\documentclass{article}

\usepackage{amsmath,amssymb}
\usepackage{xurl}
\usepackage[colorlinks=true,linkcolor=blue,citecolor=blue,urlcolor=blue]{hyperref}

% Shared commands for notes in docs/paper-gaps/.

\DeclareUnicodeCharacter{2080}{\ifmmode{}_{0}\else\textsubscript{0}\fi}
\DeclareUnicodeCharacter{2081}{\ifmmode{}_{1}\else\textsubscript{1}\fi}
\DeclareUnicodeCharacter{2082}{\ifmmode{}_{2}\else\textsubscript{2}\fi}
\DeclareUnicodeCharacter{2099}{\ifmmode{}_{n}\else\textsubscript{n}\fi}

\newcommand{\Lean}{\textsc{Lean}}
\ExplSyntaxOn
\NewDocumentCommand{\leanid}{m}{
  \ifmmode
    \textsf{\detokenize{#1}}
  \else
    \group_begin:
    \urlstyle{sf}
    \tl_set:Nn \l_tmpa_tl {#1}
    \tl_replace_all:Nnn \l_tmpa_tl {\_} {_}
    \exp_args:NV \path \l_tmpa_tl
    \group_end:
  \fi
}
\ExplSyntaxOff
\providecommand{\tr}{\operatorname{tr}}
\newcommand{\tnleanIssueBase}{https://github.com/LionSR/TNLean/issues/}
\newcommand{\tnleanPrBase}{https://github.com/LionSR/TNLean/pull/}
\newcommand{\tnleanDocBase}{https://sirui-lu.com/TNLean/docs/}
\newcommand{\ghissue}[1]{\href{\tnleanIssueBase#1}{GitHub issue \##1}}
\newcommand{\ghpr}[1]{\href{\tnleanPrBase#1}{PR \##1}}
\newcommand{\doclink}[1]{\href{\tnleanDocBase#1.html}{\path{#1}}}

% Dirac notation and the identity operator, as in blueprint/src/macros/common.tex.
% \providecommand keeps note-local definitions authoritative.
\usepackage{dsfont}
\providecommand{\ket}[1]{|#1\rangle}
\providecommand{\bra}[1]{\langle #1|}
\providecommand{\braket}[2]{\langle #1 | #2 \rangle}
\providecommand{\ketbra}[2]{|#1\rangle\!\langle #2|}
\providecommand{\Id}{\mathds{1}}
\providecommand{\one}{\mathds{1}}
\providecommand{\C}{\mathbb{C}}
\providecommand{\MN}[1]{M_{#1}(\C)}
\providecommand{\MD}{M_D(\C)}

% External referencing for paper sources.  The build helper writes sanitized
% auxiliary files under build/paper-aux/ and excludes duplicated source labels.
\usepackage{xr}

% Import a generated paper auxiliary file with a caller-supplied label prefix.
% The candidate paths support builds from docs/paper-gaps/, the repository root,
% and build/paper-aux/ itself.
\newcommand{\importpaperaux}[2]{%
  \IfFileExists{../../build/paper-aux/#2.aux}{%
    \externaldocument[#1]{../../build/paper-aux/#2}%
  }{%
    \IfFileExists{build/paper-aux/#2.aux}{%
      \externaldocument[#1]{build/paper-aux/#2}%
    }{%
      \IfFileExists{#2.aux}{%
        \externaldocument[#1]{#2}%
      }{}%
    }%
  }%
}

% General external reference macros
\newcommand{\paperref}[2]{\ref{#1-#2}}
\newcommand{\papereqref}[2]{(\ref{#1-#2})}

% CPSV16 source (1606.00608 / MPDO-22-12-17-2.tex) external document and shortcuts
\importpaperaux{cpsv16-}{1606.00608/MPDO-22-12-17-2-tnlean}
\newcommand{\cpsvref}[1]{\paperref{cpsv16}{#1}}
\newcommand{\cpsveqref}[1]{\papereqref{cpsv16}{#1}}

% Verdict marker: \gapnote{<kind>}{<status>}. Kind is the policy
% classification (clarification, local-correction, scope-restriction,
% unfaithful, false-source, open-gap); status is open, wip (elimination
% underway), resolved, or historical. Machine-read by the site index;
% prints a verdict line.
\newcommand{\gapnote}[2]{\par\noindent\textbf{Verdict:} #1 (#2).\par}


\title{Historical Note: Two-Layer Sector Refinement of CPSV16\\
\texttt{eq:II\_ABasicTensors}}
\author{}
\date{2026-05-13}

\begin{document}
\maketitle
\gapnote{scope-restriction}{historical}

\section{Status}

\textbf{Historical proof archaeology, not an active gap.}
This note records the mismatch between a retired strict-modulus specialization
and the repeated-copy BNT decomposition in CPSV16. The current
\leanid{MPSTensor.SectorDecomposition} and \leanid{MPSTensor.IsBNTCanonicalForm} surfaces retain the
raw copy weights, and
\leanid{MPSTensor.fundamentalTheorem_proportional_canonicalForm} proves the
corrected proportional theorem on BNT families with nonzero copy weights. No
adapter from the retired specialization is a current task.

\section{The source assertion}

Cirac--P\'erez-Garc\'ia--Schuch--Verstraete express a tensor $A$ in canonical
form in terms of its basis of normal tensors (BNT) $A_j$ through
Equation~\texttt{eq:II\_ABasicTensors} of
Ref.~\cite{Cirac2016MPDO_arXiv}. In the notation used here, the decomposition
is
\begin{align}
    A^i
        &=\bigoplus_{j=1}^{g}\bigoplus_{q=1}^{r_j}
          \mu_{j,q}X_{j,q}A^i_jX_{j,q}^{-1}.
    \label{eq:two-layer_general_bnt_decomposition}
\end{align}
The coefficients $\mu_{j,q}\in\C$ are arbitrary complex coefficients; the
source imposes no ordering or unit-modulus condition. The matrices $X_{j,q}$
are invertible gauge matrices, and the tensors $A_j$ are the BNT normal blocks.
The eventual linear independence of the BNT vectors
$\ket{V^{(N)}(A_j)}$ is recorded separately in the review's definition of
basis normal tensors in Ref.~\cite{Cirac2021Matrix}.

\section{The retired specialization}

An earlier formal route introduced a project-specific specialization of
(\ref{eq:two-layer_general_bnt_decomposition}). That route grouped
equal-modulus blocks into sectors, extracted a common spectral factor, and
imposed three hypotheses stronger than the source decomposition:
\begin{enumerate}
    \item \emph{Strict spectral ordering.}
    The predicate required spectral-level coefficients
    $\lambda_j\in\C$, one for each BNT sector $j$, satisfying
    \begin{align}
        |\lambda_0|
            &>|\lambda_1|>\cdots>|\lambda_{g-1}|.
        \label{eq:two-layer_strict_spectral_order}
    \end{align}
    This condition was formerly recorded as a strict-antitone property of the
    chosen spectral levels. No such ordering is present in
    (\ref{eq:two-layer_general_bnt_decomposition}).

    \item \emph{Unit-modulus within-sector factors.}
    Each weight $\mu_{j,q}$ was required to factor as
    \begin{align}
        \mu_{j,q}
            &=\lambda_j\nu_{j,q},
        \label{eq:two-layer_weight_factorization}\\
        |\nu_{j,q}|
            &=1.
        \label{eq:two-layer_unit_modulus_factor}
    \end{align}
    The general form
    (\ref{eq:two-layer_general_bnt_decomposition}) allows
    $|\mu_{j,q}|$ to vary arbitrarily within a sector. The unit-modulus
    factorization is conceptually related to the RFP characterization theorem
    at lines~543--555 of Ref.~\cite{Cirac2016MPDO_arXiv}, which concerns a
    strict subclass of the general BNT.

    \item \emph{Dominant normalization.}
    The dominant spectral level was required to satisfy
    \begin{align}
        |\lambda_0|
            &=1.
        \label{eq:two-layer_dominant_normalization}
    \end{align}
    The retired strict-modulus layer recorded this condition through the former
    field \texttt{mu\_dom\_norm\_one}. This is historical terminology, not a
    live Lean declaration, and the condition is not imposed by
    (\ref{eq:two-layer_general_bnt_decomposition}).
\end{enumerate}

\section{The surviving formal surface}

The surviving surface separates the source decomposition into two pieces.
\leanid{MPSTensor.SectorDecomposition} stores the basis blocks, their copy
multiplicities, and the raw weights $\mu_{j,q}$ through the coefficient
\begin{align}
    P.\mathtt{coeff}\,N\,j
        &=\sum_q(P.\mathtt{weight}\,j\,q)^N.
    \label{eq:two-layer_sector_power_sum_coefficient}
\end{align}
The predicate \leanid{MPSTensor.IsBNTCanonicalForm} records the BNT hypotheses used by
the fundamental-theorem route: irreducibility, left canonical normalization,
gauge-phase distinctness of equal-dimensional basis blocks, eventual linear
independence of the basis MPVs, bounded weights, and the existence of at least
one unit-modulus weight.

This surface keeps the two-layer information visible without assuming the
strict ordering (\ref{eq:two-layer_strict_spectral_order}) or the factorization
(\ref{eq:two-layer_weight_factorization}) as part of the BNT predicate. When
such a factorization is needed, it should be stated as an additional theorem
hypothesis rather than hidden in the canonical-form predicate.

\section{Historical scope restriction and its resolution}

The representative-family predicate \leanid{MPSTensor.IsNormalCanonicalFormBNT} covers
one chosen block per gauge-phase class and therefore does not itself store the
repeated-copy data $r_j>1$. The earlier route tried to bridge this restriction
through a strict-modulus adapter. That adapter was retired.

The current \leanid{MPSTensor.SectorDecomposition} surface directly permits
$P.\mathtt{copies}(j)>1$, and the SectorBNT comparison recovers repeated weights
from the power sums
\begin{align}
    \sum_q\mu_{j,q}^N.
    \label{eq:two-layer_repeated_weight_power_sum}
\end{align}
Thus the old one-copy restriction is not an outstanding obligation of the
proportional fundamental theorem. Its history is documented in
\path{docs/paper-gaps/cpsv16_ft_one_copy_scope_restriction.tex}.

\section{Formal declarations}

The relevant surviving formal declarations are:
\begin{itemize}
    \item \leanid{MPSTensor.SectorDecomposition}: the raw two-layer sector surface
    (\doclink{TNLean/MPS/SharedInfra/SectorDecomposition}).

    \item \leanid{MPSTensor.IsBNTCanonicalForm}: the BNT canonical-form predicate
    (\doclink{TNLean/MPS/FundamentalTheorem/SectorBNT/Basic}).

    \item
    \leanid{MPSTensor.IsBNTCanonicalForm.cross_overlap_basis_tendsto_zero}
    and
    \leanid{MPSTensor.IsBNTCanonicalForm.combined_family_eventually_li}
    in \doclink{TNLean/MPS/FundamentalTheorem/SectorBNT/Api}, together with
    \leanid{MPSTensor.SectorDecomposition.coeff_not_eventually_zero} in
    \doclink{TNLean/MPS/SharedInfra/SectorDecomposition}.

    \item
    \leanid{MPSTensor.fundamentalTheorem_proportional_canonicalForm}: the
    corrected proportional comparison for BNT families
    (\doclink{TNLean/MPS/FundamentalTheorem/SectorBNT/FundamentalCoord}).
\end{itemize}
See also \ghpr{1657}, on branch
\path{feat/mps-ft-path-beta-pr1-two-layer}, and the ``PR~1.5 amendment'' in
the audit memo
\path{docs/audits/2026-05-13_cpsv16_ft_path_beta_scout.md}.

\bibliographystyle{alpha}
\bibliography{references}

\end{document}
